\begin{table}[p]
\caption{\textbf{AD--CN regional effect sizes and BH-FDR-adjusted comparisons for Figure~\ref{fig:4}C--F.} Cohen's $d$ is AD minus CN divided by the pooled within-group sample s.d.}
\label{tab:regional_effect_sizes}
\centering\scriptsize\setlength{\tabcolsep}{3pt}
\begin{tabularx}{\textwidth}{llXrrrrrr}
\toprule
Panel & Tracer & ROI (ID) & $N_{AD}$ & $N_{CN}$ & Real $d$ & Syn $d$ & Real $q$ & Syn $q$ \\
\midrule
C & AV45 & L precuneus (1025) & 179 & 511 & 1.6783 & 1.6437 & $7.73\times10^{-66}$ & $1.25\times10^{-63}$ \\
C & AV45 & L inferiortemporal (1009) & 179 & 511 & 1.6831 & 1.6337 & $4.68\times10^{-66}$ & $4.74\times10^{-63}$ \\
C & AV45 & L middletemporal (1015) & 179 & 511 & 1.6482 & 1.6236 & $4.67\times10^{-64}$ & $1.72\times10^{-62}$ \\
C & AV45 & L superiortemporal (1030) & 179 & 511 & 1.5780 & 1.5434 & $5.64\times10^{-60}$ & $8.51\times10^{-58}$ \\
C & AV45 & L transversetemporal (1034) & 179 & 511 & 1.5133 & 1.5288 & $2.84\times10^{-56}$ & $5.61\times10^{-57}$ \\
D & FDG & L isthmuscingulate (1010) & 303 & 371 & -1.2531 & -1.1774 & $2.69\times10^{-48}$ & $2.10\times10^{-43}$ \\
D & FDG & R isthmuscingulate (2010) & 303 & 371 & -1.3972 & -1.3201 & $1.18\times10^{-57}$ & $1.94\times10^{-52}$ \\
D & FDG & L supramarginal (1031) & 303 & 371 & -0.4848 & -0.4308 & $2.14\times10^{-9}$ & $1.10\times10^{-7}$ \\
D & FDG & Left Hippocampus (17) & 302 & 371 & -0.6401 & -0.6042 & $4.27\times10^{-15}$ & $1.27\times10^{-13}$ \\
D & FDG & Right Hippocampus (53) & 303 & 371 & -0.7518 & -0.7189 & $4.59\times10^{-20}$ & $1.54\times10^{-18}$ \\
E & AV45 & Left VentralDC (28) & 179 & 511 & -0.1050 & -0.0933 & 0.2368 & 0.2858 \\
E & AV45 & Right VentralDC (60) & 179 & 511 & -0.1211 & -0.0986 & 0.1729 & 0.2618 \\
E & AV45 & Right Hippocampus (53) & 179 & 511 & -0.0068 & 0.0006 & 0.9373 & 0.9949 \\
E & AV45 & Left Cerebellum White Matter (7) & 179 & 511 & -0.0745 & -0.1010 & 0.3996 & 0.2529 \\
E & AV45 & Left Hippocampus (17) & 179 & 511 & 0.0701 & 0.1232 & 0.4241 & 0.1667 \\
F & FDG & Left Thalamus (10) & 303 & 371 & -0.0097 & 0.0582 & 0.9008 & 0.4667 \\
F & FDG & Right Thalamus (49) & 303 & 371 & -0.0668 & -0.0530 & 0.3964 & 0.4991 \\
F & FDG & L parsorbitalis (1019) & 303 & 371 & -0.1294 & -0.0913 & 0.1035 & 0.2506 \\
F & FDG & L lateraloccipital (1011) & 303 & 371 & -0.1239 & -0.0978 & 0.1171 & 0.2229 \\
F & FDG & R lingual (2013) & 303 & 371 & -0.0493 & -0.0435 & 0.5303 & 0.5744 \\
\bottomrule
\end{tabularx}
\par\smallskip\footnotesize
Separate two-sided Student's $t$-tests compare AD and CN for each tracer and PET type. BH-FDR correction uses all 98 ROIs within each of four tracer-by-PET-type families, rather than just the displayed regions. Positive $d$ indicates higher uptake in AD. All values are derived from min--max-normalized regional intensities, not SUV/SUVR. Preserved significance or similar effect sizes do not establish clinical equivalence.
\end{table}
